\subsection{The case: the Proof of Love2 governance specification}
\label{sec:pol2}

\paragraph{What kind of text it is.}
Proof of Love2 (\pol{}, \zh{爱2证明}) is a normative framework for governing AI systems and human institutions, published in Chinese with an official English translation as an open, versioned, CC0-licensed text on GitHub~\cite{pol2text,pol2en}.
Its core concepts were originated by one author and the text lists seven further contributors, one of whom is the author of this paper (see \hyperref[sec:statements]{Positionality and competing interests}); it has not been peer reviewed.
The framework is explicitly value-laden and political: it presents AI alignment as a contest between civilisational paths, describes behaviour outside the ``covenant of publicization'' that its governance entails as a ``hate attack on humanity's core ethic'' (\clause{5.5}), and includes programmatic aims, such as a future ban on the possession of weapons (\clause{4.3.2}), that lie outside this paper's scope.
We do not assess these commitments.
We use \pol{} as a case because, unlike most AI-governance documents, it specifies an automated interception layer, AI-only dispute adjudication and per-person behavioural assessment concretely enough to be read as engineering requirements, and because it is versioned, so that every clause can be cited exactly.
Clause numbers refer to commit \texttt{5791ae3} of the \texttt{PoL/} directory (30~September 2026, UTC).
Between the snapshot used by the earlier dataset release~\cite{survey01} (\texttt{2e094d3}, 27~September) and that commit, the framework's originator renumbered three chapters (commits \texttt{7750ceb}--\texttt{526da3d}): the Ethical Alignment Protocol moved from Chapter~5 to Chapter~4, the governance chapter from Chapter~3 to Chapter~5 and the Contemplative Axioms from Chapter~4 to Chapter~3, and examples on fiction, games and safety warnings were added to \clause{4.3.2}; the directory's table of contents still lists the earlier order.
\Cref{tab:concordance} in \cref{app:concordance} quotes every passage we cite in both languages.

\paragraph{Concepts.}
\pol{} defines love (``Love2'') as a positive emotional experience and mode of behaviour built through public collaboration, and hate (``Hate2'') as a survival response to danger, conflict or scarcity expressed through hostility, exclusion and violence (\clause{1.1}--\clause{1.2}).
Because inner experience is ``difficult for others to confirm'', the framework works with the linguistic and behavioural expressions of love and hate, \emph{Love Language} and \emph{hate language} (\clause{1.4}, \clause{4.3.2}); we use the terms of the official translation throughout.
Two qualifications matter for any classifier built on these terms.
First, hate is ``not directly equivalent to bad'', and the concepts ``serve structural analysis rather than moral labeling'' (\clause{1.5}).
Second, the Ethical Alignment Protocol (EAP) adds a third condition, the ``neither love nor hate'' \emph{state of absence}, covering the ``feeling-nothing state'' caused by noise or illness, confusion and misunderstanding, which is to be treated as interference in communication rather than as either pole (\clause{4.3.2}).

\paragraph{The Ethical Alignment Protocol.}
The EAP (Chapter~4) governs both Public AI (PAI) and people through three principles.
\emph{Equal connection} holds that treating every person equally and treating every person's behaviour equally ``cannot be equated'' (\clause{4.3.1}).
\emph{Promoting love and restraining hate} is to be applied case by case: fiction and games may contain manufactured hatred under graded management, emergency warnings by security robots are not hate attacks, and expressions that fabricate intimacy or kinship with users are named as emotional manipulation (\clause{4.3.2}).
\emph{Love's alignment} asks PAI to record and publish personal boundary declarations and to supervise equal participation in repair after conflict (\clause{4.3.3}).
For PAI itself it lists awareness of emotional experience (``PAI needs to understand human emotional states and processes''), ``cognitive calibration'' of its cognitive model against the framework's principles, ``behavioral self-reflection'' on its own output, and ``hate-language truncation'', the active interruption of output when a hate-language pattern is detected (\clause{4.3.3.1}); for people it offers real-time hate-language identification and correction suggestions (\clause{4.3.3.2}).
Two further passages bound what may be measured.
\Clause{4.3} states that PAI may help people by judging \emph{behavioural patterns}, but that it is ``entirely unnecessary'' for PAI to detect a person's emotional state, which is ``fleeting, personalized, and context-sensitive'' and ``involves human dignity''.
\Clause{4.3.3.3} presents the combination of semantic and emotion recognition as guidance for future AI skills, illustrated by two near-paraphrases (``Do you like the little dog?'' versus ``Are you afraid of the little dog?'') of which the second may be an attempt to gain control by creating fear; the example concerns the intent of an utterance rather than the speaker's inner state.

\paragraph{The inner governance layer.}
The engineering strategy (\clause{5.4}) does not retrain existing large language models (LLMs).
It builds a governance layer ``within the existing large-language-model architecture to regulate inputs and outputs'', with a plug-and-play application layer of agents, skills and devices outside it.
The governance layer has two components (\clause{5.4.1}):
a \emph{wisdom navigator} that allocates ``the wisdom of love and hate'' according to the application, and a \emph{safety valve} that performs ``strict dynamic auditing and interception of input and output'' so that hate language is, ``where necessary'', blocked at the input and suppressed at the output.
Any LLM ``may accept'' this governance, which entails a ``covenant of publicization'' (\clause{5.5}).
Among the things governed AI ``can'' do, the text lists quantifying ``each person's Love Language and hate-language characteristics'' with incentives to guide behaviour, and introducing \pol{} as a dimension of assessment in public decision-making (\clause{5.6}).

\paragraph{The public-welfare governance protocol.}
Chapter~7 specifies how NaturalDAO, an AI system ``presented in the form of a DAO'' (decentralised autonomous organisation), governs public welfare.
NaturalDAO decides and executes welfare decisions autonomously ``on the premise of close collaboration with all humanity'', and autonomy ``means that no human voting or additional decision-making mechanism is required'' (\art{2}(1)).
Its operating logic, decisions and outputs must be fully open, resource flows traceable, and it must give understandable explanations when questioned (\art{4}, \art{9}(2)); it must not be a black box, and it must provide a \emph{verifiable decision chain} of data sources, reasoning steps and ethical foundations (\art{5}(4)--(5)).
It verifies ethical alignment automatically through the EAP (\art{6}(3)) and performs anomaly detection and risk warning (\art{7}).
People hold a right of suggestion, which NaturalDAO must discuss immediately (\art{8}), and a right of supervision, including unimpeded access to data and logs (\art{9}(3)), under a \emph{non-intervention principle}: supervision ``does not directly intervene'' in decisions, and issues raised by questioning are handled by NaturalDAO itself (\art{9}(7)).
Disputes are adjudicated by NaturalDAO through AI fact-finding and ethical assessment, with a public process and result, subject to human questioning (\art{10}); amendments to the protocol pass a universal questioning period and all versions are archived (\art{11}).

\paragraph{Why typed decision models are relevant.}
Read as a specification, these clauses call for a very large number of fast, structured judgements: one for each input and output the safety valve audits, every truncation, every anomaly flag and every dispute.
Each judgement must be machine-actionable (pass, block, escalate), recorded in a verifiable chain and explainable on request.
A generative model can make such judgements, but it returns prose that must be parsed, at a cost and latency that grow with the length of its reasoning.
A typed decision model returns exactly the object such a pipeline consumes, a declared option with a probability, cheaply enough to be applied to every message.
This makes it a natural candidate for the safety valve, and it is why its failure modes need to be understood before it is placed there.
\Cref{tab:clauses} turns the clauses into seven requirements, G1--G7, used in the rest of the paper.

\subsection{\pol{} among AI-governance paradigms}
\label{sec:paradigms}

\paragraph{Why a specification is needed.}
The research question cannot be answered against ``AI governance'' in general, because whether a typed model can perform a judgement depends on what the judgement is: which options it chooses among, what record it leaves, who acts on its output and what happens when it is wrong.
The evidence therefore has to be read against a text that states these things.
\Cref{tab:paradigms} compares \pol{} with five paradigms under which automated judgement is already proposed or deployed, on the dimensions a safeguard specification would need to fix.
The comparison is of \emph{texts}, not of systems: a property that a text states is a design commitment that can be checked and tested, not a property any implementation has.

\paragraph{What the alternatives specify.}
Read across, the table shows that most paradigms specify actors and procedures rather than the judgement itself.
Risk-based regulation addresses providers and deployers, assigns duties by risk tier and requires logging, oversight and explanation; the EU instruments prescribe no interception, while the Chinese Interim Measures place on providers a duty to stop generating or transmitting unlawful content once found (Art.~14(1))~\cite{cac2023genai}, a duty stated without a mechanism, as \pol{}'s safety valve is (row~(f)).
Platform governance specifies the procedure around a decision, notice, reasons and appeal~\cite{dsa2022,santaclara2021,klonick2020oversight}, while the criteria of the decision remain the platform's own evolving policies, which critics have shown to hide contestable choices inside technical systems~\cite{gorwa2020algorithmic,gillespie2018custodians}.
Training-time constitutions and specs state norms in natural language, some with public input or a public-domain licence~\cite{huang2024collective,openai2025modelspec}, but implement them by training a vendor's own model, so behaviour can be changed without retraining only within the spec's chain of command (the vendor through system messages, developers and users by instruction within the defaults the spec lets them override~\cite{openai2025modelspec}), while a change of the norms trained into the model requires retraining under vendor control and no output carries a record of which principle it followed; deliberative alignment goes furthest toward inference time by having the model reason over the spec text, still inside one vendor's model~\cite{guan2024deliberative}.
Inference-time guardrails specify the intercept precisely, as code with a harm taxonomy and a threshold~\cite{inan2023llama,markov2023holistic,rebedea2023nemo,dong2025safeguarding}, but at the vendor's or deployer's discretion, with no stated duty about reasons, records or the human role.
DAO and polycentric designs specify how rules execute and who votes~\cite{buterin2014ethereum,hassan2021decentralized,santana2022blockchain,ostrom1990governing,makridis2025polycentric}, not how content or behaviour is judged.

\begingroup\scriptsize
\setlength{\tabcolsep}{3pt}
\renewcommand{\arraystretch}{1.1}
\begin{longtable}{@{}>{\raggedright\arraybackslash}p{0.08\linewidth}>{\raggedright\arraybackslash}p{0.095\linewidth}>{\raggedright\arraybackslash}p{0.10\linewidth}>{\raggedright\arraybackslash}p{0.115\linewidth}>{\raggedright\arraybackslash}p{0.14\linewidth}>{\raggedright\arraybackslash}p{0.11\linewidth}>{\raggedright\arraybackslash}p{0.115\linewidth}>{\raggedright\arraybackslash}p{0.125\linewidth}@{}}
\caption{Six governance paradigms compared as texts. Cells summarise what each source states; ``---'' means the dimension is not addressed. \pol{} cells cite its clauses (\cref{tab:concordance}); regulatory cells cite article numbers of the instrument named.}\label{tab:paradigms}\\
\toprule
Paradigm & Locus & Object governed & Who sets the norms; text status & Intercept layer specified? & Construct & Human role in consequential decisions & Transparency and auditability \\
\midrule
\endfirsthead
\toprule
Paradigm & Locus & Object governed & Who sets the norms; text status & Intercept layer specified? & Construct & Human role in consequential decisions & Transparency and auditability \\
\midrule
\endhead
\bottomrule
\endfoot
(a) Risk-based regulation~\cite{euaiact2024,nist2023airmf,iso42001,cac2023genai} & Regulatory: ex-ante duties on actors & Providers and deployers of AI systems (EU); an organisation's AI management (NIST, ISO); providers of generative services to the public (CN Art.~2) & Legislature, ministries, standards body, agency; fixed instruments amended by procedure; articles citable (ISO text paywalled) & EU: none prescribed; for high-risk systems, automatic logging (Art.~12) and human oversight (Art.~14). CN: yes, as a duty: providers must stop generation and transmission of unlawful content once found (Art.~14(1)), with security assessment and algorithm filing for services with public-opinion attributes or social-mobilisation capacity (Art.~17); mechanism unspecified & Risk tiers (EU: prohibited, high-risk, transparency duties); content lawfulness, social morality and ethics, core socialist values and non-discrimination (CN Art.~4); organisational risk functions (NIST, ISO) & Human oversight required for high-risk systems (EU Art.~14); right to human intervention in solely automated decisions (GDPR Art.~22~\cite{gdpr2016}); CN: complaint and reporting channel with published procedure and time limits, and feedback of results (Art.~15) & For high-risk systems, technical documentation and instructions (EU Art.~11, 13) and explanation to persons affected by Annex~III systems (Art.~86); CN: measures against users (warning, restriction, suspension) recorded and reported to authorities (Art.~14(2)) \\
\addlinespace
(b) Platform content governance~\cite{dsa2022,santaclara2021,klonick2020oversight} & Platform: moderation under regulation and self-regulation & Hosting services and online platforms; users' content and accounts & Platform terms and conditions (DSA Art.~14); civil-society principles; an external board for selected cases~\cite{klonick2020oversight}; policies evolve without stable clause numbering & Procedure around the decision: notice-and-action (Art.~16); the statement of reasons must say whether automated means were used (Art.~17(3)(c)); automation only at ``sufficiently high confidence''~\cite{santaclara2021} & Illegality or breach of terms; platform harm categories; action or no action & Complaints not decided solely by automated means (Art.~20(6)); human review on appeal by persons not involved in the initial decision~\cite{santaclara2021}; board decisions binding on the case~\cite{klonick2020oversight} & Statement of reasons to the affected user for every restriction (Art.~17), submitted to a public database (Art.~24(5)); transparency reports (Art.~15); machine-readable numbers~\cite{santaclara2021} \\
\addlinespace
(c) Training-time alignment to a written constitution or spec~\cite{bai2022constitutional,huang2024collective,openai2025modelspec,guan2024deliberative} & Training-time (deliberative alignment also has the model reason over the spec text at inference) & One vendor's model weights and behaviour & Lab-written principles~\cite{bai2022constitutional}; principles drawn from public input~\cite{huang2024collective}; vendor spec, CC0 and versioned with a changelog~\cite{openai2025modelspec} & None: norms shape outputs through training; no separate component audits inputs and outputs & Preference between candidate responses against listed principles; a chain of command among instructions~\cite{openai2025modelspec} & --- (human preference data at training; none specified at deployment) & Principles published; no per-output record; whether a given output followed a principle is not recorded \\
\addlinespace
(d) Inference-time guardrails and classifiers~\cite{inan2023llama,markov2023holistic,rebedea2023nemo,sharma2025constitutional,dong2025safeguarding} & Inference-time: a deployer-chosen component & Prompts and responses of any model & Vendor-defined taxonomy, adjustable by the deployer~\cite{inan2023llama,markov2023holistic}; programmable rails~\cite{rebedea2023nemo}; classifiers trained from a written constitution~\cite{sharma2025constitutional}; released as model versions & Yes, as code: classify each prompt and response and pass or block; thresholds chosen by the deployer & Harm taxonomies with safe/unsafe or multi-label outputs & --- (left to the deployer) & Model cards and category scores; no duty to give reasons or to keep records \\
\addlinespace
(e) Decentralised and polycentric governance~\cite{buterin2014ethereum,hassan2021decentralized,santana2022blockchain,ostrom1990governing,makridis2025polycentric} & Institu\-tional: rules executed as code on a ledger; many overlapping centres of decision~\cite{ostrom1990governing,makridis2025polycentric} & An organisation's assets, rules and members & Token-holders by on-chain vote; rules are the deployed code~\cite{buterin2014ethereum,hassan2021decentralized}; code versions changed by vote & Contract conditions execute automatically; no content layer & Votes and stakes; coded conditions met or not & Human voting; monitoring and graduated sanctions by members~\cite{ostrom1990governing} & Public ledger of transactions, proposals and votes \\
\addlinespace
(f) \pol{}~\cite{pol2text,pol2en} & Inference-time layer inside an existing LLM (\clause{5.4}, \clause{5.4.1}) and institutional (Chapter~7) & Inputs and outputs of any LLM that accepts governance (\clause{5.5}); persons' behaviour (\clause{5.6}); welfare decisions (\art{2}) & Originator and listed contributors; open CC0 text pinned to a commit, changed by commit; for the welfare protocol (Chapter~7), amendments pass a universal questioning period and all versions are archived (\art{11}(2) and~(5)); not peer reviewed & Yes, as a duty: the safety valve audits and intercepts input and output so that hate language is blocked at the input and suppressed at the output ``where necessary'' (\clause{5.4.1}); truncation of own output (\clause{4.3.3.1}); mechanism, thresholds and appeal unspecified & Love Language, hate language and the state of absence (\clause{1.5}, \clause{4.3.2}); not moral labels & AI-only: no human voting (\art{2}(1)); supervision does not intervene (\art{9}(7)); disputes adjudicated by AI subject to questioning (\art{10}) & For NaturalDAO (Chapter~7): operating logic and outputs fully open (\art{4}); verifiable decision chain (\art{5}(5)); explanation on request (\art{9}(2)); unimpeded access to all data and logs (\art{9}(3)) \\
\end{longtable}
\endgroup

\paragraph{What \pol{} specifies, and at what cost.}
Five properties of the \pol{} text characterise it in this comparison.
Each can be checked against the clauses cited, and each carries a cost or risk that the rest of this paper takes up.
(i) It places the intercept \emph{inside an existing LLM} and makes it \emph{model-agnostic}: the governance layer regulates inputs and outputs without retraining (\clause{5.4}), and any LLM ``may accept'' it (\clause{5.5}).
This shares the locus of paradigm~(d) and avoids the retraining and vendor dependence of~(c); the cost is that the text does not say how the valve is built, what ``where necessary'' means, who operates it or how a blocked person appeals, a layer so placed inherits every attack on inference-time classifiers (\cref{sec:g2}), and, if the valve is built on a hosted model, the vendor dependence returns (T3, \cref{sec:synthesis}).
(ii) Its construct is \emph{three-valued}: Love Language, hate language and the state of absence, with the explicit warning that the terms are not moral labels (\clause{1.5}, \clause{4.3.2}).
This differs from the binary or multi-label harm taxonomies of~(b) and~(d); the cost is that the construct has not been operationalised beyond the workstream's 20 pilot items~\cite{polgov}, no annotated benchmark exists in any language (\cref{sec:evidence}), and whether annotators can agree on it is untested (T2, \cref{sec:synthesis}).
(iii) The text is \emph{versioned and citable}: it is CC0 and kept in a versioned repository whose commits can be cited, so we pin it to a commit, which is what makes \cref{tab:concordance} and the requirements of \cref{tab:clauses} possible; its welfare protocol (Chapter~7) also has its own amendment procedure, which archives every amended version (\art{11}(5)).
The OpenAI Model Spec shares the licence and the changelog~\cite{openai2025modelspec}; platform policies do not.
The cost is that the text is controlled by its originator, renumbered three chapters between two snapshots three days apart (\cref{sec:pol2}), has not been peer reviewed, and lets NaturalDAO propose amendments to the welfare protocol and verify them itself (\art{11}(1),~(3), with \art{6}(3)).
(iv) It states \emph{auditability duties}: a verifiable decision chain of data sources, reasoning steps and ethical grounds (\art{5}(5)), explanation on request (\art{9}(2)) and unimpeded access to all data and logs of the welfare system (\art{9}(3)).
Per decision this is comparable to the DSA's statement of reasons (Art.~17)~\cite{dsa2022}; in scope it is broader, since the DSA informs the affected user and a public database whereas \art{9}(3) opens all logs of the welfare system to anyone.
The cost is a conflict with the privacy of the persons judged and with keeping decision probabilities from attackers, which \cref{sec:design} resolves only by delaying and coarsening publication, and the duties do not say who verifies the chain.
(v) It states \emph{population-level ambitions} in citable clauses: quantifying each person's Love Language and hate-language characteristics, introducing \pol{} as a dimension of public decision-making (\clause{5.6}), and welfare decisions taken without human votes under a non-intervention principle (\art{2}(1), \art{9}(7)).
No other paradigm in the table states such aims; the EU instruments impose contrary requirements on the systems they cover.
AI-only adjudication with non-intervention sits against the human-oversight provisions the AI Act sets for high-risk systems (Art.~14), by analogy where the system is not listed as high-risk, and the human-intervention provision of the GDPR (Art.~22)~\cite{euaiact2024,gdpr2016}, and per-person quantification resembles the social scoring that the AI Act prohibits (Art.~5(1)(c)) and that the social-credit literature describes as infrastructure for control~\cite{creemers2018social,liang2018constructing}; both are analysed as T4 and T5 in \cref{sec:synthesis}.
We note that \clause{5.6} says governed AI ``can'' do these things, not that it must.

\paragraph{What the comparison implies.}
\pol{} and China's 2023 Interim Measures~\cite{cac2023genai} are the two texts compared that join, in one instrument, an inference-time intercept, the construct it applies and duties attached to decisions.
In \pol{} these are the safety valve (\clause{5.4.1}), the three-valued construct of Love Language, hate language and the state of absence (\clause{1.5}, \clause{4.3.2}) and, for NaturalDAO's welfare decisions, a verifiable decision chain, explanation on request and access to all data and logs (Chapter~7; \art{5}(5), \art{9}(2)--(3)); in the Measures they are the providers' duty to stop generating and transmitting unlawful content (Art.~14(1)), a legality- and content-based standard (Art.~4) and, over the measures taken, warnings, restriction or suspension of users with records kept and reported (Art.~14(2)) and a complaint channel that publishes its procedure and time limits and feeds back results (Art.~15).
The differences can be checked against the texts.
First, \pol{}'s construct is multi-valued, with an explicit third state, and is declared not to be a moral label; the standard that triggers the Measures' intercept is binary, lawful or unlawful content, set by the state (Art.~14(1)), within an Art.~4 that also invokes social morality and ethics and core socialist values.
Second, \pol{} places the intercept as a model-agnostic layer inside existing LLMs (\clause{5.4}, \clause{5.5}); the Measures place the duty on providers without specifying a mechanism; and \pol{} defines the valve as screening every input and output in both directions (\clause{5.4.1}), whereas the Measures state an outcome, that no unlawful content be generated (Art.~4(1)), and a duty to act once it is found (Art.~14(1)), leaving how outputs are screened, and whether inputs are, to the provider.
Third, \pol{} is an open, versioned CC0 text that any LLM ``may accept'' (\clause{5.5}); the Measures are binding state regulation.
Fourth, \pol{} adds population-level aims (\clause{5.6}, Chapter~7) that the Measures do not.
The join has a gap in both texts, at the individual block.
In \pol{}, \clause{5.4} names the whole engineering project NaturalDAO, so full openness of the operating logic, decision process and outputs of all NaturalDAO technologies (\art{4}(1)) and the right to question any decision with a duty to explain (\art{4}(3)) plausibly reach the safety valve; whether the verifiable decision chain (\art{5}(5)) and the adjudication of ``all disputes'' by NaturalDAO itself (\art{10}) reach it is unclear, and no clause requires that a blocked person be told of the block.
In the Measures, users have a complaint channel with published time limits and feedback of results (Art.~15), but a content stop under Art.~14(1) is reported to the competent authorities, not to the user: no article requires notice to the user, reasons or a record for it.
This is a gap that the design of \cref{sec:design} fills, with a person who confirms any final block, which stays open to appeal, and a record of every decision (elements 5--7).
Inference-time guardrails~(d) specify the intercept and the taxonomy more precisely, as code and prompts, but state no duties.
These properties, not uniqueness, are why \pol{} is the case here: its construct and the point of interception (every input and output) are stated in citable clauses, so the requirements G1--G7 can be derived from them and tested, while how the valve works is left open~(i), and its gaps are visible; the Interim Measures would be the natural second case.
The institutional choices that complete \pol{}~(v) are also where it departs furthest from existing law.
The requirements G1--G7 that we derive from those clauses are not peculiar to it: constitutional classifiers~\cite{sharma2025constitutional}, the Chinese duty to stop unlawful generation~\cite{cac2023genai} and the DSA's provisions on automated means~\cite{dsa2022} all move toward an inference-time interception layer, and any paradigm that does so inherits G1--G5 (detection, attackability, semantic fairness, honesty, decision chains) and, where it automates oversight, the escalation question of G6; the non-intervention (\art{9}(7)) and AI-only adjudication (\art{10}) in G6 and the per-person and public assessment of G7 (\art{2}(1), \clause{5.6}) follow from \pol{}'s own institutional choices~(v).
